\begin{exercise}{68}
Show that for categories $\mathcal C$ and $\mathcal D$ we have $\mathrm N(\mathcal C)\star\mathrm N(\mathcal D)\cong\mathrm N(\mathcal C\star\mathcal D)$. In particular, show that there is a canonical isomorphism $\Delta^i\star\Delta^j\cong\Delta^{i+1+j}$.

\begin{proof}
By Definition~1.4.8, an $n$-simplex of
$\mathrm N(\mathcal C)\star\mathrm N(\mathcal D)$
consists of a cut $[n]=J_1\amalg J_2$ and functors
$J_1\to\mathcal C$, $J_2\to\mathcal D$;
an empty part contributes a unique choice.
Such data give a functor $[n]\to\mathcal C\star\mathcal D$:
use the two functors within their respective parts and
the unique join morphism across the cut.
All mixed composition identities hold because the
corresponding cross hom-sets are singletons.
Conversely, a functor $T\colon[n]\to\mathcal C\star\mathcal D$
cannot pass from $\mathcal D$ back to $\mathcal C$, since
there are no morphisms in that direction
(Definition~1.4.1).
Thus the objects mapped into $\mathcal C$ form an initial
segment $J_1$, the rest form $J_2$, and the restrictions
of $T$ recover exactly the original data.
These constructions are inverse bijections.

Precomposition by $\alpha\colon[k]\to[n]$ replaces the
cut by $(\alpha^{-1}J_1,\alpha^{-1}J_2)$ and the two
functors by their precompositions, exactly as in the
join structure maps (Lemma~1.4.6).
Hence the bijections form a simplicial isomorphism,
canonical and natural in both categories.
Finally, Lemma~1.1.44 and Exercise~\ref{ex:67} give
\[
 \Delta^i\star\Delta^j
 \cong\mathrm N([i])\star\mathrm N([j])
 \cong\mathrm N([i]\star[j])
 \cong\mathrm N([i+j+1])
 \cong\Delta^{i+1+j}.
\]
\end{proof}
\end{exercise}
